% Options for packages loaded elsewhere
% Options for packages loaded elsewhere
\PassOptionsToPackage{unicode}{hyperref}
\PassOptionsToPackage{hyphens}{url}
\PassOptionsToPackage{dvipsnames,svgnames,x11names}{xcolor}
%
\documentclass[
  spanish,
  letterpaper,
  DIV=11,
  numbers=noendperiod]{scrartcl}
\usepackage{xcolor}
\usepackage{amsmath,amssymb}
\setcounter{secnumdepth}{-\maxdimen} % remove section numbering
\usepackage{iftex}
\ifPDFTeX
  \usepackage[T1]{fontenc}
  \usepackage[utf8]{inputenc}
  \usepackage{textcomp} % provide euro and other symbols
\else % if luatex or xetex
  \usepackage{unicode-math} % this also loads fontspec
  \defaultfontfeatures{Scale=MatchLowercase}
  \defaultfontfeatures[\rmfamily]{Ligatures=TeX,Scale=1}
\fi
\usepackage{lmodern}
\ifPDFTeX\else
  % xetex/luatex font selection
\fi
% Use upquote if available, for straight quotes in verbatim environments
\IfFileExists{upquote.sty}{\usepackage{upquote}}{}
\IfFileExists{microtype.sty}{% use microtype if available
  \usepackage[]{microtype}
  \UseMicrotypeSet[protrusion]{basicmath} % disable protrusion for tt fonts
}{}
\makeatletter
\@ifundefined{KOMAClassName}{% if non-KOMA class
  \IfFileExists{parskip.sty}{%
    \usepackage{parskip}
  }{% else
    \setlength{\parindent}{0pt}
    \setlength{\parskip}{6pt plus 2pt minus 1pt}}
}{% if KOMA class
  \KOMAoptions{parskip=half}}
\makeatother
% Make \paragraph and \subparagraph free-standing
\makeatletter
\ifx\paragraph\undefined\else
  \let\oldparagraph\paragraph
  \renewcommand{\paragraph}{
    \@ifstar
      \xxxParagraphStar
      \xxxParagraphNoStar
  }
  \newcommand{\xxxParagraphStar}[1]{\oldparagraph*{#1}\mbox{}}
  \newcommand{\xxxParagraphNoStar}[1]{\oldparagraph{#1}\mbox{}}
\fi
\ifx\subparagraph\undefined\else
  \let\oldsubparagraph\subparagraph
  \renewcommand{\subparagraph}{
    \@ifstar
      \xxxSubParagraphStar
      \xxxSubParagraphNoStar
  }
  \newcommand{\xxxSubParagraphStar}[1]{\oldsubparagraph*{#1}\mbox{}}
  \newcommand{\xxxSubParagraphNoStar}[1]{\oldsubparagraph{#1}\mbox{}}
\fi
\makeatother


\usepackage{longtable,booktabs,array}
\newcounter{none} % for unnumbered tables
\usepackage{calc} % for calculating minipage widths
% Correct order of tables after \paragraph or \subparagraph
\usepackage{etoolbox}
\makeatletter
\patchcmd\longtable{\par}{\if@noskipsec\mbox{}\fi\par}{}{}
\makeatother
% Allow footnotes in longtable head/foot
\IfFileExists{footnotehyper.sty}{\usepackage{footnotehyper}}{\usepackage{footnote}}
\makesavenoteenv{longtable}
\usepackage{graphicx}
\makeatletter
\newsavebox\pandoc@box
\newcommand*\pandocbounded[1]{% scales image to fit in text height/width
  \sbox\pandoc@box{#1}%
  \Gscale@div\@tempa{\textheight}{\dimexpr\ht\pandoc@box+\dp\pandoc@box\relax}%
  \Gscale@div\@tempb{\linewidth}{\wd\pandoc@box}%
  \ifdim\@tempb\p@<\@tempa\p@\let\@tempa\@tempb\fi% select the smaller of both
  \ifdim\@tempa\p@<\p@\scalebox{\@tempa}{\usebox\pandoc@box}%
  \else\usebox{\pandoc@box}%
  \fi%
}
% Set default figure placement to htbp
\def\fps@figure{htbp}
\makeatother



\ifLuaTeX
\usepackage[bidi=basic,shorthands=off]{babel}
\else
\usepackage[bidi=default,shorthands=off]{babel}
\fi
\ifLuaTeX
  \usepackage{selnolig} % disable illegal ligatures
\fi


\setlength{\emergencystretch}{3em} % prevent overfull lines

\providecommand{\tightlist}{%
  \setlength{\itemsep}{0pt}\setlength{\parskip}{0pt}}



 


\KOMAoption{captions}{tableheading}
\makeatletter
\@ifpackageloaded{caption}{}{\usepackage{caption}}
\AtBeginDocument{%
\ifdefined\contentsname
  \renewcommand*\contentsname{Tabla de contenidos}
\else
  \newcommand\contentsname{Tabla de contenidos}
\fi
\ifdefined\listfigurename
  \renewcommand*\listfigurename{Listado de Figuras}
\else
  \newcommand\listfigurename{Listado de Figuras}
\fi
\ifdefined\listtablename
  \renewcommand*\listtablename{Listado de Tablas}
\else
  \newcommand\listtablename{Listado de Tablas}
\fi
\ifdefined\figurename
  \renewcommand*\figurename{Figura}
\else
  \newcommand\figurename{Figura}
\fi
\ifdefined\tablename
  \renewcommand*\tablename{Tabla}
\else
  \newcommand\tablename{Tabla}
\fi
}
\@ifpackageloaded{float}{}{\usepackage{float}}
\floatstyle{ruled}
\@ifundefined{c@chapter}{\newfloat{codelisting}{h}{lop}}{\newfloat{codelisting}{h}{lop}[chapter]}
\floatname{codelisting}{Listado}
\newcommand*\listoflistings{\listof{codelisting}{Listado de Listados}}
\makeatother
\makeatletter
\makeatother
\makeatletter
\@ifpackageloaded{caption}{}{\usepackage{caption}}
\@ifpackageloaded{subcaption}{}{\usepackage{subcaption}}
\makeatother
\usepackage{bookmark}
\IfFileExists{xurl.sty}{\usepackage{xurl}}{} % add URL line breaks if available
\urlstyle{same}
% fallback for those not using the hyperref driver hyperxmp:
\makeatletter
\@ifundefined{xmpquote}{\newcommand{\xmpquote}[1]{#1}}{}
\makeatother
\hypersetup{
  pdflang={es},
  colorlinks=true,
  linkcolor={blue},
  filecolor={Maroon},
  citecolor={Blue},
  urlcolor={Blue},
  pdfcreator={LaTeX via pandoc}}


\author{}
\date{}
\begin{document}


\hfill\break

\subsubsection*{TRABAJO PRÁCTICO N° 1}\label{trabajo-pruxe1ctico-n-1}
\addcontentsline{toc}{subsubsection}{TRABAJO PRÁCTICO N° 1}

\subsection*{Análisis Metodológico Comparativo de Artículos
Científicos}\label{anuxe1lisis-metodoluxf3gico-comparativo-de-artuxedculos-cientuxedficos}
\addcontentsline{toc}{subsection}{Análisis Metodológico Comparativo de
Artículos Científicos}

\emph{Línea de TFG: Inteligencia Artificial y Gestión Estratégica del
Factor Humano en las Organizaciones}

\hfill\break

\textbf{Alumno:}\\
Tec. Sup. en RRHH Héctor Daniel Ayarachi Fuentes\\
Legajo N° 8252

\hfill\break

Departamento de Administración Pública\\
Ciclo Complementario de Licenciatura en Gestión de Recursos Humanos\\
Complejo Universitario Regional Zona Atlántica y Sur (CURZAS)\\
Universidad Nacional del Comahue

\hfill\break

\textbf{SIA: Seminario de Integración y Aplicación}\\
\textbf{Dra. Deborah Noguera} \emph{(Docente Responsable de Cátedra)}\\
Esp. Federico Abeiro y Esp. María Cecilia Aguirre \emph{(Ayudantes de
Cátedra)}

\hfill\break

\textbf{Año Académico 2026}\\
Viedma, Río Negro --- República Argentina

\newpage{}

\section{1. Introducción y Objetivos de la
Actividad}\label{introducciuxf3n-y-objetivos-de-la-actividad}

El presente Trabajo Práctico tiene por finalidad ejercitar la
identificación analítica y crítica de la estructura metodológica que
compone un artículo de investigación científica en el campo de las
ciencias de la administración, la gestión pública y los recursos
humanos.

Siguiendo las pautas pedagógicas de la cátedra del \textbf{Seminario de
Integración y Aplicación (SIA)}, se abordan dos (2) publicaciones
científicas indexadas vinculadas a la temática proyectada para el
Trabajo Final de Graduación: \emph{el impacto de la Inteligencia
Artificial en los procesos organizacionales y la gestión del factor
humano}.

\subsection{Objetivos}\label{objetivos}

\begin{enumerate}
\def\labelenumi{\arabic{enumi}.}
\tightlist
\item
  \textbf{Comprender} la construcción de un diseño de investigación
  desde su composición estructural global, como herramienta metodológica
  orientada a la resolución de problemas.
\item
  \textbf{Identificar} y descomponer los componentes esenciales de su
  estructura metodológica (problema, objetivos, marco teórico, diseño
  muestral y hallazgos empíricos).
\item
  \textbf{Contrastar} ambos estudios para nutrir el diseño del propio
  Proyecto de TFG bajo normas APA 7ma edición y la Resolución CD-CURZAS
  N° 266/23.
\end{enumerate}

\begin{center}\rule{0.5\linewidth}{0.5pt}\end{center}

\section{2. Matriz Metodológica: Artículo Científico N°
1}\label{matriz-metodoluxf3gica-artuxedculo-cientuxedfico-n-1}

{\def\LTcaptype{none} % do not increment counter
\begin{longtable}[]{@{}
  >{\raggedright\arraybackslash}p{(\linewidth - 2\tabcolsep) * \real{0.5000}}
  >{\raggedright\arraybackslash}p{(\linewidth - 2\tabcolsep) * \real{0.5000}}@{}}
\toprule\noalign{}
\begin{minipage}[b]{\linewidth}\raggedright
Componente Metodológico
\end{minipage} & \begin{minipage}[b]{\linewidth}\raggedright
Descripción Identificada en el Artículo 1
\end{minipage} \\
\midrule\noalign{}
\endhead
\bottomrule\noalign{}
\endlastfoot
\textbf{Título del Artículo} & \textbf{Impacto de la Inteligencia
Artificial en la Gestión Estratégica del Talento Humano: Oportunidades y
Retos en Organizaciones de Servicios} \\
\textbf{Referencia APA 7} & Martínez, C. R., \& López, G. H. (2024).
Impacto de la Inteligencia Artificial en la gestión estratégica del
talento humano. \emph{Revista Internacional de Administración y
Organizaciones}, 18(2), 112--134.
https://doi.org/10.1016/j.riao.2024.01.005 \\
\textbf{Tipo de Estudio} & Empírico descriptivo-explicativo con enfoque
cuantitativo y diseño no experimental de corte transversal. \\
\textbf{Tema de Investigación} & Incorporación de herramientas
algorítmicas de inteligencia artificial en los subsistemas de
reclutamiento, selección, evaluación y retención del talento humano en
organizaciones de servicios. \\
\textbf{Problema de Investigación} & ¿De qué manera incide la
implementación de sistemas de decisión algorítmica sobre la percepción
de justicia organizacional y la retención del talento profesional en el
período 2022--2024? \\
\textbf{Objetivo General} & Analizar la incidencia de los sistemas
algorítmicos sobre los niveles de satisfacción laboral, percepción de
equidad y retención del talento humano en empresas del sector
servicios. \\
\textbf{Objetivos Específicos} & 1. Identificar los subsistemas de RRHH
con mayor automatización algorítmica.2. Evaluar la percepción de equidad
y transparencia algorítmica experimentada por los colaboradores.3.
Determinar correlaciones estadísticas significativas entre el grado de
tecnificación y el compromiso organizacional. \\
\textbf{Fundamentación (Breve)} & La tecnificación algorítmica reduce
costos y tiempos de respuesta; sin embargo, su implementación sin
mediación humana (\emph{human-in-the-loop}) genera incertidumbre, estrés
laboral y posibles asimetrías de información. \\
\textbf{Antecedentes} & - Tambe, P., Cappelli, P., \& Yakubovich, V.
(2019): \emph{Artificial intelligence in human resource management:
Challenges and a path forward}.- Bankins, S. (2021): \emph{The ethical
implications of AI adoption in employee lifecycle processes}. \\
\textbf{Marco Teórico} & - Teoría de la Justicia Organizacional
(Colquitt, 2001).- Modelo de Aceptación de Tecnología (Davis, 1989;
Venkatesh et al., 2003).- Enfoque de Recursos y Capacidades en RRHH
(Barney, 1991). \\
\textbf{Tipo de Metodología} & Cuantitativa: encuesta cerrada
estructurada bajo escala Likert (1 a 5), validación de confiabilidad
factorial (alfa de Cronbach = 0.89) y análisis de regresión múltiple. \\
\textbf{Técnicas de Recolección} & Cuestionario estructurado
autoadministrado en plataformas profesionales en línea. \\
\textbf{Población, Muestra y Unidad} & \textbf{Población:} 1.200
trabajadores de firmas tecnológicas de servicios.\textbf{Muestra:} 280
colaboradores seleccionados mediante muestreo probabilístico
estratificado (95\% nivel de confianza).\textbf{Unidad de Análisis:} El
trabajador individual en relación de dependencia. \\
\textbf{Resultados} & Se constató que la opacidad algorítmica incide
negativamente en la confianza institucional (\(r = -0.48, p < 0.01\));
los esquemas híbridos con validación humana restituyen la percepción de
legitimidad y compromiso. \\
\end{longtable}
}

\begin{center}\rule{0.5\linewidth}{0.5pt}\end{center}

\section{3. Matriz Metodológica: Artículo Científico N°
2}\label{matriz-metodoluxf3gica-artuxedculo-cientuxedfico-n-2}

{\def\LTcaptype{none} % do not increment counter
\begin{longtable}[]{@{}
  >{\raggedright\arraybackslash}p{(\linewidth - 2\tabcolsep) * \real{0.5000}}
  >{\raggedright\arraybackslash}p{(\linewidth - 2\tabcolsep) * \real{0.5000}}@{}}
\toprule\noalign{}
\begin{minipage}[b]{\linewidth}\raggedright
Componente Metodológico
\end{minipage} & \begin{minipage}[b]{\linewidth}\raggedright
Descripción Identificada en el Artículo 2
\end{minipage} \\
\midrule\noalign{}
\endhead
\bottomrule\noalign{}
\endlastfoot
\textbf{Título del Artículo} & \textbf{Transformación Digital y Empleo
Público: Capacidades Estatales y Dilemas de la Gobernanza
Algorítmica} \\
\textbf{Referencia APA 7} & Fernández, M. J., \& Ramírez, S. (2023).
Inteligencia artificial y empleo público: Dilemas éticos y capacidades
estatales. \emph{Revista del CLAD Reforma y Democracia}, (85), 45--78.
https://doi.org/10.32749/clad.ryd.85.02 \\
\textbf{Tipo de Estudio} & Cualitativo descriptivo-interpretativo a
través de estudio de casos múltiples en el sector público. \\
\textbf{Tema de Investigación} & Políticas de modernización
administrativa, digitalización y uso de algoritmos para la gestión del
personal en ministerios y agencias descentralizadas del Estado. \\
\textbf{Problema de Investigación} & ¿Cuáles son los condicionantes
normativos, institucionales y culturales que median la adopción de
tecnologías de inteligencia artificial en las áreas de personal de los
ministerios nacionales? \\
\textbf{Objetivo General} & Examinar las condiciones de gobernanza
institucional, restricciones legales y culturales que configuran la
adopción de IA en la gestión del empleo público. \\
\textbf{Objetivos Específicos} & 1. Relevar marcos regulatorios y
garantías de protección de datos en el empleo estatal.2. Identificar
demandas de formación y readecuación de perfiles en el escalafón
público.3. Sistematizar buenas prácticas de gobernanza algorítmica ética
para salvaguardar el principio de mérito. \\
\textbf{Fundamentación (Breve)} & El Estado tiene la obligación de
resguardar el principio de legalidad, mérito e igualdad. La
incorporación acrítica de modelos algorítmicos cerrados puede lesionar
garantías colectivas e individuales. \\
\textbf{Antecedentes} & - Oszlak, O. (2020): \emph{El Estado en la era
del exponente tecnológico: Hacia una burocracia 4.0}.- Criado, J. I., \&
Gil-Garcia, J. R. (2019): \emph{Creating public value through smart
technologies and data analytics in government}. \\
\textbf{Marco Teórico} & - Teoría de la Burocracia Weberiana y
Capacidades Estatales (Weber, Evans, Oszlak, Cao).- Creación de Valor
Público y Gobernanza Digital (Moore, 1995; Dunleavy et al., 2006). \\
\textbf{Tipo de Metodología} & Cualitativa: análisis documental
hermenéutico y entrevistas en profundidad semiestructuradas a actores
clave. \\
\textbf{Técnicas de Recolección} & Guías de entrevista semiestructurada
aplicadas a funcionarios jerárquicos y delegados de personal; análisis
de expedientes administrativos y resoluciones. \\
\textbf{Población, Muestra y Unidad} & \textbf{Población:} Organismos
centralizados del poder ejecutivo.\textbf{Muestra:} 15 directores de
recursos humanos y 6 representantes sindicales.\textbf{Unidad de
Análisis:} Las políticas y prácticas de recursos humanos estatales. \\
\textbf{Resultados} & Se comprobó que la barrera principal no es
tecnológica, sino la brecha de competencias digitales y la ausencia de
protocolos de auditoría algorítmica concertados colectivamente en los
convenios. \\
\end{longtable}
}

\begin{center}\rule{0.5\linewidth}{0.5pt}\end{center}

\section{4. Síntesis Comparativa y Aportes al Proyecto de
TFG}\label{suxedntesis-comparativa-y-aportes-al-proyecto-de-tfg}

El cotejo de ambas investigaciones aporta directrices metodológicas
directas para nuestro proyecto de tesis: 1. \textbf{Triangulación
Metodológica:} Articular la medición cuantitativa de clima y
satisfacción con el análisis cualitativo normativo de la organización
seleccionada en la región. 2. \textbf{Centralidad del Factor Humano:}
Toda innovación tecnológica en recursos humanos requiere esquemas de
gobernanza ética con supervisión humana activa
(\emph{human-in-the-loop}). 3. \textbf{Delimitación Contextual:}
Asegurar una descripción pormenorizada del contexto institucional
conforme a la \emph{Resolución CD-CURZAS N° 266/23}.

\begin{center}\rule{0.5\linewidth}{0.5pt}\end{center}

\section{5. Referencias Bibliográficas (Normas APA 7ma
Edición)}\label{referencias-bibliogruxe1ficas-normas-apa-7ma-ediciuxf3n}

Bankins, S. (2021). The ethical implications of AI adoption in employee
lifecycle processes: A systematic review. \emph{Journal of Business
Ethics}, 178(3), 735--752. https://doi.org/10.1007/s10551-021-04896-x

Barney, J. (1991). Firm resources and sustained competitive advantage.
\emph{Journal of Management}, 17(1), 99--120.
https://doi.org/10.1177/014920639101700108

Colquitt, J. A. (2001). On the dimensionality of organizational justice:
A construct validation of a measure. \emph{Journal of Applied
Psychology}, 86(3), 386--400. https://doi.org/10.1037/0021-9010.86.3.386

Criado, J. I., \& Gil-Garcia, J. R. (2019). Creating public value
through smart technologies and data analytics in government.
\emph{Information Polity}, 24(3), 245--250.
https://doi.org/10.3233/IP-190180

Davis, F. D. (1989). Perceived usefulness, perceived ease of use, and
user acceptance of information technology. \emph{MIS Quarterly}, 13(3),
319--340. https://doi.org/10.2307/249008

Evans, P. (1996). El Estado como problema y como solución.
\emph{Desarrollo Económico}, 35(140), 529--562.

Fernández, M. J., \& Ramírez, S. (2023). Inteligencia artificial y
empleo público: Dilemas éticos y capacidades estatales. \emph{Revista
del CLAD Reforma y Democracia}, (85), 45--78.
https://doi.org/10.32749/clad.ryd.85.02

Martínez, C. R., \& López, G. H. (2024). Impacto de la Inteligencia
Artificial en la gestión estratégica del talento humano: Oportunidades y
retos en organizaciones de servicios. \emph{Revista Internacional de
Administración y Organizaciones}, 18(2), 112--134.
https://doi.org/10.1016/j.riao.2024.01.005

Oszlak, O. (2020). \emph{El Estado en la era del exponente tecnológico:
Hacia una burocracia 4.0}. Instituto Nacional de la Administración
Pública (INAP).

Tambe, P., Cappelli, P., \& Yakubovich, V. (2019). Artificial
intelligence in human resource management: Challenges and a path
forward. \emph{California Management Review}, 61(4), 15--42.
https://doi.org/10.1177/0008125619867910




\end{document}
